\documentclass[12pt,a4paper]{article}
\usepackage{circuitikz}
\usepackage{amsmath}
\usetikzlibrary{circuits.logic.US,calc}


\begin{document}
\begin{center}
    \begin{circuitikz}[american,thick, x=1.3cm ]

    % \draw (0, 0) node[op amp] (op1) {};
    % \draw (5,-0.5) node[op amp] (op2) {};
    % \draw (op1.+) to[R , l=$3k\Omega$] (-4,-0.5) to[V , l=$10V$] (-4, -4);
    % \draw (-4,-4) -- ( 7,-4);
    % \draw (op1.-) -- (-1,0.5) -- (-1, 2) -- (1,2)-- ( 1,0);
    % \draw (-1.5,-0.5) to[R , l_=$10k\Omega$] (-1.5,-4); 
    % \draw (-1, -0.5) -- (-1,-2) to[R, l=$6k\Omega$] ( 3,-2) -- (3,2) to[R, l=$5k\Omega$]  (7,2)  --(7,-0.5) --(op2.out);
    %  \draw (7, -0.5) to[R, l=$5k\Omega$, i=$i_o$] (7,-4);
    % \draw (op1.out) to[R, l=$5k\Omega$] (3,0) -- (op2.-);
    % \draw (op2.+) -- (4, -1) -- (4,-4);
    % \draw (3,-4) node[ground]{};

    \draw (0,0) node[op amp] (op1){};
    \draw (op1.-) to[R ,l=$5k\Omega$] (-3,0.5) to[V , l=$2v$]  (-3,-3);
    \draw (op1.+) to (-1, -0.5) to[R, l=$10k\Omega$] (-1,-3) -- (-3,-3);
    \draw (-2,-3) node[ground]{};
    \draw (op1.-) -- (-0.9,1.5)  to[R, l=$20k\Omega$] (1.5,1.5)--(1.5,0);
    \draw (op1.out)  to[R, l=$40k\Omega$] (4,0);

    

    \draw (0,-8) node[op amp] (op2){};
    \draw (op2.-) to[R ,l=$10k\Omega$] (-3,-7.5) to[V , l=$3v$]  (-3,-13);
    \draw (op2.+) to (-1, -8.5) to[R, l=$30k\Omega$] (-1,-13) -- (-3,-13);
    \draw (-2,-13) node[ground]{};
    \draw (op2.out)  to[R, l=$80k\Omega$] (4,-8) to[R, l=$20k\Omega$] (4,-11) node[ground]{};
    \draw (-1, -10) to[R, l=$50k\Omega$] (1.5,-10)--(1.5,-8);


    \draw (5,-4) node[op amp] (op3){};
    \draw (op3.-) -- (4,0);
    \draw (op3.+) -- (4,-8);

    \draw (8.5,-4) node[op amp, rotate=180](op4){};

    \draw (op3.out)  to[R, l=$20k\Omega$, i=$i_o$] (op4.out);
    \draw (op4.-) to[R, l_=$20k\Omega$] (9.3,-7) to[V , l=$4v$] (9.3,-10) node[ground]{};


    \draw (4,0) to[R, l=$40k\Omega$] (6,0)-- (op3.out);
    \draw (op4.out) -- (7.5,-2)-- (9.5,-2)-- (op4.+);
\end{circuitikz}
\end{center}
\newpage




\begin{center}
\begin{circuitikz}[thick]
    % Set style to IEEE
    % \ctikzset{logic ports=us}
    
    % Define gates
    \node[and port] (AND1) at (0,0) {};
    \node[not port] (NOT1) at (0,-3) {};
    \node[not port] (NOT2) at (0,-7) {};
    \draw (AND1.in 1) -- ++(-4,0) node[left]{$P$};
    \draw (AND1.in 2) -- ++(-4,0) node[left]{$Q$};

    \draw (NOT2.in) -- ++ (-3,0) --++(0,7.3);

    \node[or port, anchor=in 2] (OR1) at (1,-6.5){};
    \draw (NOT2.out) -- (OR1.in 2);
    \draw (OR1.in 1) -- ++ (-4,0) --++ (0,5.7);
    \draw (NOT1.in) -- ++(-2.3,0);

    \node[and port] (AND2) at (3.5,-4){};
    \draw (AND2.in 1) -- (NOT1.out);
    \draw (AND2.in 2) --++ (-5.8,0);

    \node[or port] (OR2) at (6,-1){};
    \draw (AND1.out) -- (OR2.in 1);
    \draw (AND2.out) -- (OR2.in 2);

    \node[and port] (AND3) at (9, -2.5){};
    \draw (AND3.in 2) -- (OR1.out);
    \draw (AND3.in 1) -- (OR2.out);
    \draw (AND3.out) --++ (0.5,0);



    
    % Connect inputs and outputs
    % \draw (OR1.out) -- (AND1.in 1);
    % \draw (NOT1.out) -- (AND1.in 2);
    
    % % Label inputs/outputs
    % \draw (OR1.in 1) -- ++(-1,0) node[left]{$A$};
    % \draw (OR1.in 2) -- ++(-1,0) node[left]{$B$};
    % \draw (NOT1.in) -- ++(-1,0) node[left]{$C$};
    % \draw (AND1.out) -- ++(1,0) node[right]{$Out$};
\end{circuitikz}
\end{center}

\begin{center}
\begin{circuitikz}[thick, scale=0.9, transform shape]

    
\end{circuitikz}
\end{center}


\begin{center}
        \begin{circuitikz}[american, thick, y=1.6cm]
        
        % Transistor (with circle)
        \draw (0,0) node[npn] (Q) {};
        \draw (Q.B) -- (-2,0) to[battery,l=$V_B$] (-2,-2) node[ground]{};

        \draw (Q.C) -- (0,1) -- (1,1) node[ocirc] {} node[right] {$V_o$};
        \draw (0,1) to[R] (0,3);
        \draw (Q.E) to[R, l=$R_E$] (0,-2)  node[ground]{};
        \draw[->] (0, 3) -- (0, 3.2) node[above] {$15\text{ V}$};


        
        % % Base Connection (V3 and 22k ohm resistor)
        % \draw (Q.B) -- (-1.5, 0) coordinate(b_junct);
        % \draw (-4, 0) node[ground] {} to[R, l=$22\text{ k}\Omega$] (b_junct);
        % \draw (b_junct) to[short, *-o] (-1.5, 1) node[above] {V3};
        
        % % Collector Connection (V1 and 1.6k ohm resistor)
        % \draw (Q.C) -- (0, 0.5) coordinate(c_junct);
        % \draw (c_junct) to[R, l=$1.6\text{ k}\Omega$] (0, 2.5) coordinate(vcc);
        % \draw (c_junct) to[short, *-o] (1.5, 0.5) node[right] {V1};
        
        % % +5V Power Terminal (Hollow Triangle pointing up)
        % \draw (vcc) -- (-0.2, 2.5) -- (0, 2.8) -- (0.2, 2.5) -- cycle;
        % \node[left] at (-0.1, 2.7) {$+5\text{V}$}; 
        
        % % Emitter Connection (V2 and 2.2k ohm resistor)
        % \draw (Q.E) -- (0, -0.5) coordinate(e_junct);
        % \draw (e_junct) to[R, l_=$2.2\text{ k}\Omega$] (0, -2.5) coordinate(vee);
        % \draw (e_junct) to[short, *-o] (1.5, -0.5) node[right] {V2};
        
        % % -5V Power Terminal (Hollow Triangle pointing down)
        % \draw (vee) -- (-0.2, -2.5) -- (0, -2.8) -- (0.2, -2.5) -- cycle;
        % \node[left] at (-0.1, -2.7) {$-5\text{V}$}; 
        
    \end{circuitikz}
\end{center}
\newpage
\begin{center}
    \begin{circuitikz}[thick, american]
        \node[op amp] (O1) at (0,0) {};
        \node[op amp] (O2) at (-4.5,-0.5){};
        \draw (O2.out) to[R, l=$R$] (O1.+);
        \node[op amp] (O3) at (-9,-1){};
        \draw (O3.out) to[R, l=$R$] (O2.+);
        \draw(O3.-) -- ++(-1,0)coordinate(o3) --++(0,1.5)to[R, l=$R_f$]  ++(3.4,0) to[short,-*] (O3.out);
        \draw(O2.-) -- ++(-1,0) --++(0,1.5) -- ++(3.4,0) to[short,-*] (O2.out);
        \draw(O1.-) -- ++(-1,0) --++(0,1.5) -- ++(3.4,0) to[short,-*] (O1.out);
        \draw(O3.+) --++(0,-3)--++(13.4,0) node[ground, midway]{};
        \draw(O2.+) to[C,l=$C$]++(0,-3.5);
        \draw(O1.+) to[C,l=$C$]++(0,-4);
        \draw(O1.out) to[R,l=$R$] ++(2,0) coordinate(o1);
        \draw(o1) to[C,l=$C$, v=$V_f$]++(0,-4.5);
        \draw(o1) to[short, *-] ++(0,3.5)--++(-16,0) coordinate(j);
        \draw(o3) --++(-1.6,0) to[R,l=$R_i$] (j);

        \draw(O3.out) to[open,v=$V_o$] ++(0,-3.5);


        \draw (O2.up) -- ++(0,0.4) node[above] {$+15$ V};
        \draw (O1.up) -- ++(0,0.4) node[above] {$+15$ V};
        \draw (O3.up) -- ++(0,0.4) node[above] {$+15$ V};
        \draw (O3.down) -- ++(0,-0.4) node[below] {$-15$ V};
        \draw (O2.down) -- ++(0,-0.4) node[below] {$-15$ V};
        \draw (O1.down) -- ++(0,-0.4) node[below] {$-15$ V};





        






        
    \end{circuitikz}
\end{center}


\end{document}